\chapter{LEXICAL ANALYSIS AND SYMBOL TABLE}

This chapter covers the responsibilities related to lexical and early semantic
foundations of the language. It includes defining token structures and
categories (such as keywords, operators, identifiers, and literals), implementing
the lexer to transform source characters into tokens, and establishing the
symbol table to manage variable and function storage along with their
associated types. Together, these components ensure that raw source input is
systematically organised into meaningful structures that can be used reliably
by later stages of the compiler.

\section{Token representation}
The token system, written in \texttt{components/tokens.py}, establishes the
vocabulary and structure of all lexical elements in the language. It specifies
both the types of tokens that can exist and the structure each token must
follow. This serves as a canonical contract between the lexer and the parser,
ensuring consistency across all stages of the compiler.

Table~\ref{tab:token-categories} summarises the main categories, example lexemes,
and typical \texttt{TokenType} values.

\begin{table}[H]
\centering
\caption{Token categories and representative \texttt{TokenType} values}
\label{tab:token-categories}
\small
\begin{tabular}{p{0.16\textwidth}p{0.32\textwidth}p{0.45\textwidth}}
\toprule
Category & Example lexemes & Representative \texttt{TokenType} \\
\midrule
Literals &
\texttt{42}, \texttt{3.14}, \texttt{true}, \texttt{"hello"} &
\texttt{INT\_LITERAL}, \texttt{FLOAT\_LITERAL}, \texttt{BOOL\_LITERAL},
\texttt{STRING\_LITERAL} \\
Identifiers &
\texttt{x}, \texttt{counter}, \texttt{add} &
\texttt{IDENTIFIER} \\
Keywords &
\texttt{if}, \texttt{else}, \texttt{while}, \texttt{def},
\texttt{return}, \texttt{print} &
keyword-specific token variants \\
Operators &
\texttt{+}, \texttt{-}, \texttt{*}, \texttt{/}, \texttt{=},
\texttt{==}, \texttt{!=} &
\texttt{PLUS}, \texttt{MINUS}, \texttt{STAR}, \texttt{SLASH},
\texttt{ASSIGN}, \texttt{EQUAL\_EQUAL}, \texttt{BANG\_EQUAL} \\
Delimiters &
Parentheses, braces, comma, and semicolon (see lexer delimiter tokens) &
delimiter token variants \\
End marker & end of source input & \texttt{EOF} \\
\bottomrule
\end{tabular}
\end{table}

Token types (from class \texttt{TokenType}) are defined using an \texttt{Enum}
with \texttt{auto()} to automatically assign unique values. This avoids reliance
on raw strings and ensures safer comparisons during parsing. The token
categories include literals, identifiers, keywords, operators, delimiters, and
a special end-of-file (\texttt{EOF}) token. The \texttt{EOF} token is appended
after all input is processed, allowing the parser to detect the end of input
explicitly rather than encountering an unexpected termination.

Each token is represented as an immutable dataclass (\texttt{frozen=True}),
meaning its fields cannot be modified after creation. Each token contains four
key pieces of information: its kind (\texttt{token\_type}), its lexeme (the
exact substring from the source code), and its source position (\texttt{line},
\texttt{column}). This structured representation allows later stages to
process the program in terms of syntax and errors without handling raw text
directly. Immutability improves reliability by preventing accidental
modification during later processing.

\begin{lstlisting}
@dataclass(frozen=True)
class Token:
   token_type: TokenType
   lexeme: str
   line: int
   column: int
\end{lstlisting}

\section{Lexer implementation}
The lexer, implemented in \texttt{components/lexica.py}, is a hand-written
scanner that converts raw source text into a sequence of tokens. Unlike
generated lexers (for example SLY lexer mode, or Lex and Flex), this
implementation uses explicit control flow (\texttt{if} and \texttt{while}) to
process input.

The lexer maintains internal state including the source string, current
position, start of the current token, line number, and column index. The
column index is tracked per line and resets upon encountering a newline, while
a separate \texttt{token\_column} records the starting position of each token
for accurate error reporting.

The main function, \texttt{tokenize}, iterates through the source code, repeatedly
invoking \texttt{\_scan\_token} to identify the next token. Tokens are appended
to a list, and an \texttt{EOF} token is added at the end. Whitespace is skipped,
while invalid input results in a \texttt{LexerError}.

\begin{lstlisting}
def tokenize(self) -> list[Token]:
    tokens: list[Token] = []
    while not self._is_at_end():
        self.start = self.current
        self.token_column = self.column
        token = self._scan_token()
        if token is not None:
            tokens.append(token)
    tokens.append(Token(TokenType.EOF, "", self.line, self.column))
    return tokens
\end{lstlisting}

Token recognition follows a structured approach. Single-character tokens are
mapped using a lookup table. Multi-character operators such as \texttt{==} and
\texttt{!=} are handled using lookahead via \texttt{\_match}. String literals
are processed by \texttt{\_string}, which reads characters until a closing
quotation mark or the end of input is reached, with error handling for
unterminated strings. Numeric literals are handled by \texttt{\_number}, which
distinguishes integers and floating-point values based on the presence of a
decimal point followed by digits. Identifiers are processed greedily, consuming
alphanumeric characters and underscores, and then checked against a keyword map
to determine whether they are reserved words or user-defined names.

\begin{lstlisting}
KEYWORDS: dict[str, TokenType] = {
   "if": TokenType.IF,
   "else": TokenType.ELSE,
   "while": TokenType.WHILE,
   # ...
}
\end{lstlisting}

\begin{lstlisting}
SINGLE_CHAR_TOKENS: dict[str, TokenType] = {
   "+": TokenType.PLUS,
   "-": TokenType.MINUS,
   "*": TokenType.STAR,
   "/": TokenType.SLASH,
   # ...
}
\end{lstlisting}

Helper functions such as \texttt{\_advance}, \texttt{\_peek}, and
\texttt{\_peek\_next} manage character consumption and lookahead, while
\texttt{\_is\_at\_end} determines when input has been fully processed. Error
messages include precise line and column information, aiding debugging and
diagnostics.

Table~\ref{tab:lexica-functions} lists the main functions and data structures
in \texttt{components/lexica.py} and their roles.

\begin{table}[H]
\centering
\caption{Key functions and data in \texttt{lexica.py}}
\label{tab:lexica-functions}
\small
\begin{tabular}{p{0.25\textwidth}p{0.4\textwidth}p{0.28\textwidth}}
\toprule
Function / data & Responsibility & Output / effect \\
\midrule
\texttt{tokenize()} & Iterates through source, scans tokens, appends
\texttt{EOF} & ordered token list \\
\texttt{\_scan\_token()} & Dispatches recognition based on character and
lookahead & next token or skip \\
\texttt{\_number()} & Parses numeric sequences, distinguishes int vs float &
numeric literal token \\
\texttt{\_string()} & Parses quoted text, checks termination & string token or
error \\
\texttt{\_identifier()} & Parses identifiers, resolves keyword vs user-defined
& keyword or identifier token \\
\texttt{KEYWORDS} map & Maps reserved words to token types & token type
selection \\
\texttt{SINGLE\_CHAR\_TOKENS} map & Maps single-character symbols & token type
selection \\
\texttt{\_advance()}, \texttt{\_peek()}, \texttt{\_match()} & Character
consumption and lookahead & scanner state update \\
\bottomrule
\end{tabular}
\end{table}

\section{Symbol table and semantic structure}
The symbol table, defined in \texttt{components/symbol\_table.py}, operates at a
higher level than the lexer. While the lexer processes raw text into tokens, the
symbol table manages semantic information such as variable and function names,
their types, and their scope.

The symbol table is implemented as a scoped structure, where each scope
maintains its own dictionary of symbols and optionally references a parent
scope. This enables proper handling of nested blocks and variable shadowing.
Symbols are introduced through definition functions and retrieved through
lookup operations, which search recursively through parent scopes if needed.

\begin{lstlisting}
def lookup(self, name: str) -> Symbol:
    symbol = self._symbols.get(name)
    if symbol is not None:
        return symbol
    if self.parent is not None:
        return self.parent.lookup(name)
    raise SymbolTableError(f"Undefined symbol: {name}")
\end{lstlisting}

The type system is represented using a \texttt{LanguageType} enumeration, which
includes fundamental types such as Integer, Float, Boolean, String, and Void.
Symbols are modelled using an inheritance hierarchy: the base \texttt{Symbol}
class stores a name and type, while subclasses such as
\texttt{VariableSymbol} and \texttt{FunctionSymbol} include additional
attributes. Variables track initialization state, while functions store
parameter lists and return types. Table~\ref{tab:symbol-table-components}
summarises the main building blocks.

\begin{table}[H]
\centering
\caption{Symbol table components and roles}
\label{tab:symbol-table-components}
\small
\begin{tabular}{p{0.24\textwidth}p{0.32\textwidth}p{0.36\textwidth}}
\toprule
Component & Stored information & Role in semantic analysis \\
\midrule
Scoped symbol table & Symbols in current and parent scope & Supports nesting
and shadowing via recursive lookup \\
\texttt{LanguageType} enum & Integer, Float, Boolean, String, Void & Defines
the type system \\
\texttt{Symbol} (base) & Name and declared or inferred type & Base abstraction
for entries \\
\texttt{VariableSymbol} & Variable metadata (including initialization state) &
Tracks correct usage before assignment \\
\texttt{FunctionSymbol} & Parameters and return type & Validates calls and
return consistency \\
\texttt{lookup(name)} & Recursive scope search & Resolves identifiers or
reports undefined \\
\texttt{SymbolTableError} & Semantic error details & Reports duplicates,
undefined symbols, invalid operations \\
\bottomrule
\end{tabular}
\end{table}

Errors related to semantic violations are handled via
\texttt{SymbolTableError}. These include duplicate declarations within the same
scope, references to undefined symbols, and invalid operations such as
attempting to mark a non-variable symbol as initialized.

\section{Role of the symbol table in the pipeline}
The symbol table is constructed during the type-checking phase rather than
during lexical analysis. This guideline is followed in the implementation.
The type checker initialises a global symbol table and populates it while
traversing the AST, including defining variables and functions, managing nested
scopes, and validating type correctness.

\begin{lstlisting}
def run_pipeline(source_code: str) -> PipelineResult:
   tokens = Lexer(source_code).tokenize()
   tree = Parser(tokens).parse()
   ast_text = ASTPrinter().print(tree)
   checked_scope = TypeChecker().check(tree)
   outputs: list[str] = []
   Interpreter(output_callback=outputs.append).execute(tree)
   return PipelineResult(
       tokens=tokens,
       ast_text=ast_text,
       checked_scope=checked_scope,
       outputs=outputs,
   )
\end{lstlisting}

Once type checking is complete, the populated symbol table is returned as part
of the pipeline result (\texttt{checked\_scope}). The checker creates child
scopes for blocks and functions during analysis, but the displayed
\texttt{checked\_scope} table represents the global scope entries returned by
the pipeline. It can then be displayed or used for further analysis.

This separation of concerns keeps lexical analysis focused on tokenization,
while semantic meaning is enforced during type checking. For user-defined
functions, definitions are first registered in the global scope with placeholder
parameter and return types. These are inferred from the first valid call and
refined through body checking; subsequent calls must match the inferred
parameter types.
